% Part: proof-theory
% Chapter: cut-elimination
% Section: ce-largest

\documentclass[../../../include/open-logic-section]{subfiles}

\begin{document}

\olfileid{pt}{cut}{inv}

\olsection{حذف بزرگ‌ترین برش‌ها}

اثبات \olref[seq][inv]{lem:G3c-invert} نشان داد که اگر \Log{G3c}
نتیجهٔ یک قاعدهٔ منطقی، جز \RightR{\lexists} و
\LeftR{\lforall}، را !!{prove} کند، هر یک از مقدمات آن قاعده را نیز
!!{prove} می‌کند، بی‌آنکه !!{height} برهان افزایش یابد. می‌توانیم نتیجه‌ای
قوی‌تر نشان دهیم: این امر برای $\Log{G3c}+\CutCS$ نیز برقرار است.

مانند پیش، \emph{رتبهٔ برش} یک استنتاج \CutCS{} را عمق
!!{formula} برش آن می‌گیریم.

\begin{lem}\ollabel{lem:inv-G3c-cut}
اگر $\Log{G3c}+\CutCS$ نتیجهٔ یک قاعدهٔ منطقی~$R$، جز
\RightR{\lexists} یا \LeftR{\lforall}، را !!{prove} کند و این کار را با
!!a{proof}~$\pi$ انجام دهد، هر مقدمه را نیز !!{prove} می‌کند، با
!!{proof}~$\pi'$ (یا با !!{proof}های~$\pi'$ و~$\pi''$، برحسب اینکه
قاعده یک یا دو مقدمه دارد). افزون بر این:
(الف)~$\pheight{\pi'}\le\pheight{\pi}$ (و
$\pheight{\pi''}\le\pheight{\pi}$)؛ و
(ب)~شمار استنتاج‌های \CutCS{} و نیز رتبه‌هایشان در~$\pi'$ (و
$\pi''$) همان است که در~$\pi$ بود.
\end{lem}

\begin{proof}
اثبات‌های
\cref{pt:seq:inv:lem:G3c-invert,pt:seq:inv:lem:invert-quant}
را بررسی کنید. آن اثبات با استقرا بر~$\pheight{\pi}$ بود. در حالت
پایه، یک اصل را با اصلی دیگر جایگزین کردیم؛ پس شرط‌های (الف) و~(ب)
بدیهی‌اند. اگر آخرین قاعدهٔ~$\pi$ همان~$R$ باشد، زیر-!!{proof}های
بی‌واسطهٔ آن همان !!{proof}های~$\pi_i$ مطلوب‌اند و هر دو شرط برقرارند.
در همهٔ حالت‌های دیگر، اگر آخرین استنتاج~$S$ باشد، فرض استقرا را بر
زیر-!!{proof}های بی‌واسطهٔ~$\pi$، یعنی برهان‌های منتهی به مقدمه یا
مقدمات~$S$، اعمال می‌کنیم و سپس همان قاعدهٔ~$S$ را دوباره بر نتایج
اعمال می‌کنیم. شرط‌های (الف) و~(ب) برای زیر-!!{proof}های بی‌واسطهٔ
برهان تازه برقرارند و بنابراین برای کل برهان نیز برقرارند. فقط باید
حالتی را بررسی کنیم که آخرین استنتاج \CutCS است. باز تنها یک نمونه را
می‌آوریم: فرض کنید $\pi$ !!a{proof}ی برای نتیجهٔ
$!B \land !C, \Gamma \Sequent \Delta$ از \LeftR{\land} باشد و به
~\CutCS پایان یابد:
\[
\AxiomC{}
\RightLabel{$\pi_1$}
\Deduce$!B \land !C, \Gamma \fCenter \Delta, !A$
\AxiomC{}
\RightLabel{$\pi_2$}
\Deduce$!A, !B \land !C, \Gamma \fCenter \Delta$
\RightLabel{\CutCS}
\BinaryInf$!B \land !C, \Gamma \fCenter \Delta$
\DisplayProof
\]
فرض استقرا بر $\pi_1$ و~$\pi_2$ اعمال می‌شود، زیرا آن‌ها نیز
!!{proof}هایی از نمونه‌های نتیجهٔ \LeftR{\land} هستند، و
!!{proof}های~$\pi_1'$ و~$\pi_2'$ را به‌ترتیب برای
$!B,!C,\Gamma \fCenter \Delta,!A$ و
$!A,!B,!C,\Gamma \fCenter \Delta$ به دست می‌دهد. آن‌ها را با
\CutCS{} ترکیب می‌کنیم تا~$\pi'$ حاصل شود:
\[
\AxiomC{}
\RightLabel{$\pi_1'$}
\Deduce$!B, !C, \Gamma \fCenter \Delta, !A$
\AxiomC{}
\RightLabel{$\pi_2'$}
\Deduce$!A, !B, !C, \Gamma \fCenter \Delta$
\RightLabel{\CutCS}
\BinaryInf$!B, !C, \Gamma \fCenter \Delta$
\DisplayProof
\]
آشکار است که (الف) و~(ب) برقرارند.
\end{proof}

توجه کنید اگر به جای \CutCS از \Cut{} استفاده می‌کردیم این استدلال
کار نمی‌کرد، زیرا نتیجهٔ آخرین !!{proof} در مقدم دو نسخه از
$!B,!C,\Gamma$ و در تالی دو نسخه از~$\Delta$ داشت. ناچار بودیم به
پذیرفتنی‌بودن انقباض استناد کنیم، اما اثبات
\olref[seq][inv]{prop:G3c-cont-adm} از لم وارون‌پذیری استفاده می‌کند؛
پس استدلال دوری می‌شد.

می‌توانیم با \cref{pt:cut:inv:lem:inv-G3c-cut} اثبات دیگری برای حذف
برش ارائه کنیم. در این اثبات رخدادهای بالاترین استنتاج‌های \CutCS{}
را پی‌درپی حذف نمی‌کنیم، بلکه رخدادهای استنتاج‌های \CutCS{} با رتبهٔ
بیشینه را حذف می‌کنیم. یعنی با !!a{proof}~$\pi$ در
$\Log{G3c}+\CutCS$ آغاز می‌کنیم، برشی را در هر جای آن حذف می‌کنیم
(و شاید آن را با برش‌هایی با رتبهٔ کمتر جایگزین کنیم)، و این کار را
تا وقتی هیچ برشی نماند ادامه می‌دهیم. تنها شرط این است که بالای برشی
که با آن کار می‌کنیم، هیچ برشی با همان رتبه یا رتبه‌ای بالاتر نباشد.

\begin{lem}\ollabel{lem:max-cut-red-G3c}
اگر $\pi$ !!a{proof}ی در $\Log{G3c}+\CutCS$ باشد که به استنتاج
\CutCS{} با رتبهٔ~$n$ پایان می‌یابد و در جای دیگر فقط از قواعد
~\Log{G3c} و استنتاج‌های \CutCS{} با رتبهٔ~$<n$ استفاده کند، آنگاه
برهانی~$\pi'$ برای همان دنبالهٔ پایانی در
$\Log{G3c}+\CutCS$ وجود دارد که هر استنتاج \CutCS{} در آن رتبهٔ
~$<n$ دارد.
\end{lem}

\begin{proof}
!!{proof}~$\pi$ چنین پایان می‌یابد:
\[
\AxiomC{}
\RightLabel{$\pi_1$}
\Deduce$\Gamma \fCenter \Delta, !A$
\AxiomC{}
\RightLabel{$\pi_2$}
\Deduce$!A, \Gamma \fCenter \Delta$
\RightLabel{\CutCS}
\BinaryInf$\Gamma \fCenter \Delta$
\DisplayProof
\]
حالت‌ها را برحسب شکل~$!A$ جدا می‌کنیم.

فرض کنید $!A$ اتمی است. چون همهٔ استنتاج‌های \CutCS{} در~$\pi_1$ و
$\pi_2$ باید رتبه‌ای کمتر از $\depth{!A}=0$ داشته باشند، هیچ استنتاج
\CutCS{}ای در آن‌ها وجود ندارد. با استقرا بر~$\pheight{\pi_2}$ نشان
می‌دهیم $\Gamma \Sequent \Delta$ برهانی بدون برش دارد. اگر
$\pheight{\pi_2}=0$، آنگاه $!A,\Gamma\Sequent\Delta$ یک اصل است. اگر
$!A$ اصلی نباشد، یک $!C$ اتمی !!a{element} هم~$\Gamma$ و هم~$\Delta$
است و خود $\Gamma\Sequent\Delta$ یک اصل است. اگر $!A$ اصلی باشد،
آنگاه $\Delta=\Delta',!A$. در این حالت $\pi_1$ به
$\Gamma\Sequent\Delta',!A,!A$ پایان می‌یابد. با اعمال
\olref[seq][inv]{prop:G3c-cont-adm}، یعنی پذیرفتنی‌بودن انقباض، برهانی
بدون برش برای $\Gamma\Sequent\Delta',!A$ به دست می‌آوریم.

اگر $\pheight{\pi_2}>0$، برهان به قاعده‌ای~$R$ پایان می‌یابد. یک
مقدمهٔ آن قاعده را در نظر بگیرید؛ شکل آن
$!A,\Gamma',\Pi\Sequent\Delta',\Lambda$ است، که $\Pi$ و~$\Lambda$
!!{formula}های فعال~$R$ هستند. اعمال
\olref[seq][adm]{prop:weak-G3c-adm} بر~$\pi_1$ و~$\pi_2$،
!!{proof}هایی برای
$\Gamma,\Gamma',\Pi\Sequent\Delta,\Delta',\Lambda,!A$ و
$!A,\Gamma,\Gamma',\Pi\Sequent\Delta,\Delta',\Lambda$ به دست می‌دهد
که فرض استقرا بر آن‌ها اعمال می‌شود. برای هر مقدمهٔ~$R$،
!!a{proof}ی برای
$\Gamma,\Gamma',\Pi\Sequent\Delta,\Delta',\Lambda$ حاصل می‌شود. با
اعمال~$R$ به $\Gamma,\Gamma\Sequent\Delta,\Delta$ می‌رسیم و
\olref[seq][inv]{prop:G3c-cont-adm} برهان بدون برش
$\Gamma\Sequent\Delta$ را فراهم می‌کند.

اگر $!A$ اتمی نباشد، حالت‌ها را برحسب !!{main operator} آن جدا
می‌کنیم. در هر حالت \olref{lem:inv-G3c-cut} را به کار می‌بریم تا
!!{proof}های مقدمات قواعد چپ و راست را که $!A$ در آن‌ها فرمول اصلی
است به دست آوریم. سپس مانند حالت~E در اثبات
\olref[top]{lem:cut-adm-G3c} پیش می‌رویم.
\begin{enumerate}
  \item $!A\ident\lnot !B$. !!{proof}های~$\pi_1$ و~$\pi_2$ به
  $\Gamma\Sequent\Delta,\lnot !B$ و
  $\lnot !B,\Gamma\Sequent\Delta$ پایان می‌یابند. بنا بر
  \olref{lem:inv-G3c-cut}، !!{proof}های~$\pi_1'$ و~$\pi_2'$ برای
  $!B,\Gamma\Sequent\Delta$ و $\Gamma\Sequent\Delta,!B$ وجود دارند.
  !!{proof}~$\pi'$ چنین است:
  \[
  \AxiomC{}
  \RightLabel{$\pi_2'$}
  \Deduce$\Gamma \fCenter \Delta, !B$
  \AxiomC{}
  \RightLabel{$\pi_1'$}
  \Deduce$!B, \Gamma \fCenter \Delta$
  \RightLabel{\CutCS}
  \BinaryInf$\Gamma \fCenter \Delta$
  \DisplayProof
  \]

  \item $!A\ident(!B\lor !C)$. !!{proof}های~$\pi_1$ و~$\pi_2$ به
  $\Gamma\Sequent\Delta,!B\lor !C$ و
  $!B\lor !C,\Gamma\Sequent\Delta$ پایان می‌یابند. بنا بر
  \olref{lem:inv-G3c-cut}، با وارون‌سازی \RightR{\lor} در~$\pi_1$،
  !!a{proof}~$\pi_1'$ برای $\Gamma\Sequent\Delta,!B,!C$ وجود دارد.
  بنا بر \olref{lem:inv-G3c-cut}، با وارون‌سازی \LeftR{\lor} در~$\pi_2$،
  !!a{proof}~$\pi_2'$ برای
  $!B,\Gamma\Sequent\Delta$ و !!a{proof}~$\pi_2''$ برای
  $!C,\Gamma\Sequent\Delta$ وجود دارند. با اعمال
  \olref[seq][adm]{prop:weak-G3c-adm} بر~$\pi_2''$ نیز
  !!a{proof}~$\pi_2'''$ برای $!C,\Gamma\Sequent\Delta,!B$ به دست
  می‌آوریم. !!{proof}~$\pi'$ چنین است:
  \[
  \AxiomC{}
  \RightLabel{$\pi_1'$}
  \Deduce$\Gamma \fCenter \Delta, !B, !C$
  \AxiomC{}
  \RightLabel{$\pi_2'''$}
  \Deduce$!C, \Gamma \fCenter \Delta, !B$
  \RightLabel{\CutCS}
  \BinaryInf$\Gamma \fCenter \Delta, !B$
  \AxiomC{}
  \RightLabel{$\pi_2'$}
  \Deduce$!B, \Gamma \fCenter \Delta$
  \RightLabel{\CutCS}
  \BinaryInf$\Gamma \fCenter \Delta$
  \DisplayProof
  \]

  \item $!A\ident(!B\land !C)$. تمرین.

  \item $!A\ident(!B\lif !C)$. !!{proof}های~$\pi_1$ و~$\pi_2$ به
  $\Gamma\Sequent\Delta,!B\lif !C$ و
  $!B\lif !C,\Gamma\Sequent\Delta$ پایان می‌یابند. بنا بر
  \olref{lem:inv-G3c-cut}، با وارون‌سازی \RightR{\lif} در~$\pi_1$،
  !!a{proof}~$\pi_1'$ برای $!B,\Gamma\Sequent\Delta,!C$ وجود دارد.
  بنا بر \olref{lem:inv-G3c-cut}، با وارون‌سازی \LeftR{\lif} در~$\pi_2$،
  !!a{proof}~$\pi_2'$ برای
  $\Gamma\Sequent\Delta,!B$ و !!a{proof}~$\pi_2''$ برای
  $!C,\Gamma\Sequent\Delta$ وجود دارند. با اعمال
  \olref[seq][adm]{prop:weak-G3c-adm} بر~$\pi_2''$ نیز
  !!a{proof}~$\pi_2'''$ برای $!C,!B,\Gamma\Sequent\Delta$ به دست
  می‌آوریم. !!{proof}~$\pi'$ چنین است:
  \[
  \AxiomC{}
  \RightLabel{$\pi_2'$}
  \Deduce$\Gamma \fCenter \Delta, !B$
  \AxiomC{}
  \RightLabel{$\pi_1'$}
  \Deduce$!B, \Gamma \fCenter \Delta, !C$
  \AxiomC{}
  \RightLabel{$\pi_2'''$}
  \Deduce$!C, !B, \Gamma \fCenter \Delta$
  \RightLabel{\CutCS}
  \BinaryInf$!B, \Gamma \fCenter \Delta$
  \RightLabel{\CutCS}
  \BinaryInf$\Gamma \fCenter \Delta$
  \DisplayProof
  \]

  \item $!A\ident\lforall[x][!B(x)]$. !!{proof}های~$\pi_1$ و~$\pi_2$
  به $\Gamma\Sequent\Delta,\lforall[x][!B(x)]$ و
  $\lforall[x][!B(x)],\Gamma\Sequent\Delta$ پایان می‌یابند. بنا بر
  \olref{lem:inv-G3c-cut}، !!a{proof}~$\pi_1'$ برای
  $\Gamma\Sequent\Delta,!B(c)$ وجود دارد.

  اکنون !!{proof}~$\pi_2$ را در نظر بگیرید. این برهان به
  $\lforall[x][!B(x)],\Gamma\Sequent\Delta$ پایان می‌یابد. از دنبالهٔ
  پایانی~$\pi_2$ رو به بالا حرکت کنید و هر رخداد
  $\lforall[x][!B(x)]$ در سمت چپ دنباله‌ای را علامت بزنید که به رخداد
  $\lforall[x][!B(x)]$ در دنبالهٔ پایانی~$\pi_2$ «می‌انجامد»:
  (الف)~$\lforall[x][!B(x)]$ را در دنبالهٔ پایانی~$\pi_2$ علامت
  بزنید. (ب)~اگر $\lforall[x][!B(x)]$ در بافت نتیجهٔ قاعده‌ای ظاهر
  شده و علامت خورده است، رخداد متناظر آن را در مقدمه یا مقدمات قاعده
  علامت بزنید. (ج)~اگر $\lforall[x][!B(x)]$ در نتیجهٔ یک قاعدهٔ
  \LeftR{\lforall} اصلی و علامت‌خورده است، رخدادهای فعال متناظر
  $\lforall[x][!B(x)]$ و~$!B(t)$ را در مقدمه علامت بزنید.

  همهٔ رخدادهای علامت‌خوردهٔ $\lforall[x][!B(x)]$ را در نظر بگیرید.
  هیچ‌یک در قاعده‌ای جز \LeftR{\lforall} فعال نیستند، زیرا فقط
  !!{formula}های فعال \LeftR{\lforall} علامت می‌خورند. همهٔ رخدادهای
  علامت‌خوردهٔ $\lforall[x][!B(x)]$ را از~$\pi_2$ حذف کنید.
  اگر حاصل یک !!{proof} درست باشد، رخدادهای علامت‌خورده هرگز فعال
  نبوده‌اند و همگی مستقیماً از اصول آمده‌اند. !!{proof} حاصل به
  ~$\Gamma\Sequent\Delta$ پایان می‌یابد و می‌توانیم $\pi$ را با آن
  جایگزین کنیم. بدین‌سان یک برش با رتبهٔ بیشینه حذف کرده‌ایم.

  در غیر این صورت، حاصل برهان درستی نیست، زیرا !!{formula}های
  $\lforall[x][!B(x)]$ را حذف کرده‌ایم که در برخی استنتاج‌های
  \LeftR{\lforall} فعال بوده‌اند. آن‌ها را به «استنتاج»هایی به شکل
  زیر تبدیل کرده‌ایم:
  \[
  \AxiomC{}
  \RightLabel{$\pi_t$}
  \Deduce$!B(t), \Pi \fCenter \Lambda$
  \RightLabel{\LeftR{\lforall}}
  \UnaryInf$\Pi \fCenter \Lambda$
  \DisplayProof
  \]
  هر بالاترین «استنتاج» از این نوع را با برهان زیر جایگزین کنید:
  \[
  \AxiomC{}
  \RightLabel{$\Subst{\pi_1'}{t}{c}[\Pi \Sequent \Lambda]$}
  \Deduce$\Gamma, \Pi \fCenter \Delta, \Lambda, !B(t)$
  \AxiomC{}
  \RightLabel{$\pi_t[\Gamma \Sequent \Delta]$}
  \Deduce$!B(t), \Gamma, \Pi \fCenter \Delta, \Lambda$
  \RightLabel{\CutCS}
  \BinaryInf$\Gamma, \Pi \fCenter \Delta, \Lambda$
  \DisplayProof
  \]
  و به سمت چپ هر دنبالهٔ پایین آن~$\Gamma$ و به سمت راستش~$\Delta$
  را بیفزایید. این کار را تکرار کنید تا همهٔ «استنتاج»های
  \LeftR{\lforall} که یک $!B(t)$ علامت‌خورده در مقدمه دارند با
  استنتاج \CutCS{} بر یک~$!B(t)$ جایگزین شوند. دنبالهٔ پایانی
  !!{proof} حاصل، که اکنون واقعاً یک !!{proof} درست است، به شکل
  $\Gamma,\dots,\Gamma\Sequent\Delta,\dots,\Delta$ است. پذیرفتنی‌بودن
  انقباض را اعمال کنید تا !!a{proof}ی برای~$\Gamma\Sequent\Delta$
  به دست آید و $\pi$ را با آن جایگزین کنید. یک برش بر
  $\lforall[x][!B(x)]$ را با شاید چندین برش بر~$!B(t)$ جایگزین
  کرده‌ایم، اما همهٔ آن‌ها رتبهٔ کمتری دارند. برش‌های موجود در~$\pi_1$
  یا~$\pi_2$ باقی می‌مانند و آن‌ها نیز همگی رتبهٔ کمتری دارند.
\end{enumerate}
\end{proof}

\begin{prob}
حالت $!A\ident(!B\land !C)$ را در اثبات
\olref{lem:max-cut-red-G3c} بررسی کنید. یعنی نشان دهید اگر
!!a{proof}~$\pi$ای وجود داشته باشد که چنین پایان می‌یابد:
\[
\AxiomC{}
\RightLabel{$\pi_1$}
\Deduce$\Gamma \fCenter \Delta, !B \land !C$
\AxiomC{}
\RightLabel{$\pi_2$}
\Deduce$!B \land !C, \Gamma \fCenter \Delta$
\RightLabel{\CutCS}
\BinaryInf$\Gamma \fCenter \Delta$
\DisplayProof
\]
و $\pi_1$ و~$\pi_2$ فقط برش‌هایی با رتبهٔ
$<\depth{!B\land !C}$ داشته باشند، آنگاه !!a{proof}~$\pi'$ای برای
$\Gamma\Sequent\Delta$ وجود دارد که فقط برش‌هایی با رتبهٔ
$<\depth{!B\land !C}$ دارد.
\end{prob}

% \begin{prob}
% Verify the case where $!A \ident \lexists[x][!B(x)]$ in the proof of
% \olref{lem:inv-G3c-cut}.
% \end{prob}

\olref{lem:max-cut-red-G3c} ثابت می‌کند که می‌توان یک استنتاج \CutCS{}
با بالاترین رتبه در !!a{proof} را با استنتاج‌های \CutCS{} با رتبهٔ
کمتر جایگزین کرد. باز می‌توانیم از این لم استفاده کنیم تا هر تعداد
استنتاج \CutCS{} را از !!a{proof} حذف کنیم: آن را پی‌درپی بر بالاترین
زیر-!!{proof}هایی اعمال می‌کنیم که به \CutCS{} بیشینه پایان می‌یابند.

\begin{thm}\ollabel{thm:cut-elim-G3c-largest}
هر !!{proof}~$\pi$ در $\Log{G3c}+\CutCS$ را می‌توان به برهانی در
~\Log{G3c} تبدیل کرد.
\end{thm}

\begin{proof}
نخست ثابت می‌کنیم همهٔ برش‌های با رتبهٔ بیشینه در~$\pi$ را می‌توان
حذف کرد. بگذارید $d$ رتبهٔ بیشینهٔ برش‌های موجود در~$\pi$ و $n$ شمار
برش‌های رتبهٔ~$d$ باشد. با استقرا بر~$n$ پیش می‌رویم. اگر $n=0$ چیزی
برای اثبات نیست. اگر $n>0$، یک بالاترین برش با رتبهٔ~$d$ را برگزینید
و بگذارید $\pi_1$ زیر-!!{proof} منتهی به آن باشد. بنا بر
\olref{lem:max-cut-red-G3c}، !!a{proof}~$\pi_1'$ای برای همان دنبالهٔ
پایانی وجود دارد که همهٔ برش‌هایش رتبهٔ~$<d$ دارند. $\pi_1$ را در
~$\pi$ با~$\pi_1'$ جایگزین کنید. !!a{proof} حاصل $n-1$ برش با
رتبهٔ~$d$ دارد و فرض استقرا اعمال می‌شود.

اکنون نتیجه با استقرا بر~$d$ به دست می‌آید. اگر $d=0$، همهٔ برش‌ها
اتمی‌اند و بنا بر نتیجهٔ پیشین می‌توان همه را حذف کرد. اگر $d>0$،
نتیجهٔ پیشین برهانی می‌دهد که در آن همهٔ برش‌های رتبهٔ~$d$ با
برش‌هایی با رتبهٔ~$<d$ جایگزین شده‌اند. چون $d$ رتبهٔ بیشینهٔ
برش‌های~$\pi$ بود، اکنون برهانی داریم که رتبهٔ بیشینهٔ برش‌های آن
$<d$ است و فرض استقرا اعمال می‌شود.
\end{proof}

\end{document}
